\subsection{初始结构的例子}

I. 结构的逆像。当 I 是由单个元素组成的集合时，关于 \((\mathbf{A}, \mathcal{G}, f)\) 称为在 f 下结构 G 的原像（当它存在时）。

* 拓扑在任何映射下总有原像 \(f\) ；但对于序结构或代数结构而言，情况并非如此。*

\begin{enumerate}
\def\labelenumi{\Roman{enumi}.}
\setcounter{enumi}{1}
\tightlist
\item
  诱导结构。设 A 是一个赋予某种结构的集合， \(\mathcal{S}\) 的物种 \(\Sigma\) 设 B 是 A 的一个子集，并设 \(j\) 为B到A的典范嵌入。则在 \(j\) 下的原像 \(\mathcal{S}\) （若存在）称为由 \(\mathcal{S}\) 在 B 上。
\end{enumerate}

序结构在其定义集合的每个子集上诱导出同种结构；但对于有向集的结构，情况并非如此。* 拓扑在其定义集合的每个子集上诱导出拓扑，但紧致拓扑一般不会诱导出紧致拓扑。集合A上的代数结构一般不会在任意子集B上诱导出同种结构；如果A上给定的结构由处处有定义的合成律组成，那么B必须对这些律中的每一条都保持稳定，但这个必要条件并不总是充分的。*

一般准则CST10为我们提供了关于诱导结构的如下传递性准则：

CST11。设B是A的子集，C是B的子集，且S是某物种的结构。 \(\Sigma\) 在A上，该结构诱导了一个结构 \(S'\) 同一物种在B上。则S诱导出一个物种结构 \(\Sigma\) 在 C 上当且仅当 \(S'\) 诱导出一个物种结构 \(\Sigma\) 在C上，以及由S和 \(S'\) 则完全相同。

CST12。设 A、A' 为两个赋予结构的集合 \(\mathcal{G}\) , \(\mathcal{S}'\) 的物种 \(\Sigma\) 。设 B 为 A 的子集，B' 为 A' 的子集。假设 \(\mathcal{G}\) （分别） \(\mathcal{G}'\) ) 诱导出物种结构 \(\Sigma\) 在 B（相应地，B'）上。如果 \(f\) 是从A到A'的一个态射，使得 \(f(B) \subset B'\) ，则映射 \(g\) 从 B 到 B' 中，与 \(f\) 在B上是一个态射（关于由 \(\mathcal{G}\) 并且 \(\mathcal{S}'\) ).

让 \(j\) （分别） \(j'\) ）设B（分别地，B'）到A（分别地，A'）的典范嵌入。根据定义，我们有 \(f \circ j = j' \circ g\) 。由于 \(f\) 并且 \(j\) 是态射，所以 \(f \circ j\) 由 \((\mathbf{MO}_{\mathbf{II}})\) 也是；但接着， \(j' \circ g\) 是态射，则映射 \(g\) 根据初始结构的定义是一个态射。

\begin{enumerate}
\def\labelenumi{\Roman{enumi}.}
\setcounter{enumi}{2}
\tightlist
\item
  乘积结构。设 \((\mathbf{A}_t)_{t \in \mathbf{I}}\) 是一族集合，并且在每个集合 \(\mathbf{A}_t\) 让 \(\mathcal{G}_t\) 是物种的一个结构 \(\Sigma\) . 设 \(\mathbf{E} = \prod_{t \in \mathbf{I}} \mathbf{A}_t\) 为这一族的乘积 \((\mathbf{A}_t)_{t \in \mathbf{I}}\) （第二章，第5节），并令 \(\mathbf{pr}_t\) 表示 \(\mathbf{E}\) 到上 \(\mathbf{A}_t\) 的投影。关于该族的初始结构（如果存在） \((\mathbf{A}_t, \mathcal{G}_t, \mathbf{pr}_t)_{t \in \mathbf{I}}\) 称为结构的积 \(\mathcal{G}_t\) .
\end{enumerate}

一族序结构总有积结构，但一族全序结构却并非总是如此。* 一族群结构总有积结构，但一族除环结构却未必如此。一族拓扑总有积结构，但一族局部紧空间结构却并非总是如此；在这种情况下，若族是有限的，则存在同种的积结构，但若族是无限的，则未必存在（参见《一般拓扑学》第一章第9节第7号命题14）。*

准则CST10给出如下关于积结构的结合性准则：

CST13. 设 \((\mathbf{A}_{\iota})_{\iota \in \mathbf{I}}\) 是一个集合族，并且对每个指标 \(\iota \in \mathbf{I}\) 让 \(\mathcal{S}_{\iota}\) 是物种 \(\Sigma\) 在 \(\mathbf{A}_{\iota}\) 上的一个结构. 设 \((\mathrm{J}_{\lambda})_{\lambda \in \mathbf{L}}\) 是……的一个划分 \(\mathbf{I}\) 。假设在每个部分积上 \(\mathbf{B}_{\lambda} = \prod_{\iota \in \mathrm{J}_{\lambda}} \mathbf{A}_{\iota}\) 族 \((\mathcal{S}_{\iota})_{\iota \in \mathrm{J}_{\lambda}}\) 具有一个积结构 \(\mathcal{S}_{\lambda}'\) 。那么族 \((\mathcal{S}_{\iota})_{\iota \in \mathbf{I}}\) 具有一个积结构 \(\mathcal{S}\) 当且仅当族 \((\mathcal{S}_{\lambda}')_{\lambda \in \mathbf{L}}\) 具有一个积结构 \(\mathcal{S}'\) ，并且 \(\mathbf{E} = \prod_{\iota \in \mathbf{I}} \mathbf{A}_{\iota}\) ，赋予 \(\mathcal{S}\) ，映到 \(\mathbf{F} = \prod_{\lambda \in \mathbf{L}} \mathbf{B}_{\lambda}\) ，赋予 \(\mathcal{S}'\) 的典范映射（第二章，§5，第5节）是一个同构。

CST10的另一个应用给出了关于由积结构诱导的结构的如下判据：

CST14. 设 \((\mathbf{A}_{t})_{t\in I}\) 是一个集合族，并且对每个 \(i\in I\) 让 \(\mathcal{S}_i\) 是物种 \(\Sigma\) 在 \(\mathbf{A}_i\) 上的一个结构；对于每一个 \(i\in I\) ，让 \(\mathbf{B}_i\) 为 \(\mathbf{A}_i\) 的一个子集。假设每个 \(\mathcal{S}_i\) 诱导一个结构 \(\mathcal{S}_i'\) 在 \(\mathbf{B}_i\) 上，并且在乘积 \(\mathbf{E} = \prod_{i\in I}\mathbf{A}_i\) 上存在一个结构 \(\mathcal{S}_0\) ，它是族 \((\mathcal{S}_i)\) 的积，那么以下陈述是等价的：

\begin{enumerate}
\def\labelenumi{(\alph{enumi})}
\item
  在集合上 \(\mathbf{B} = \prod_{i\in \mathbf{I}}\mathbf{B}_i\subset \mathbf{E}\) 存在一个结构 \(\mathcal{S}\) 由...诱导的 \(\mathcal{S}_0\)
\item
  的积，在集合B上存在一个结构 \(\mathcal{S}'\) 是结构族 \((\mathcal{S}_t')\) .
\end{enumerate}

此外，这些陈述蕴含 \(\mathcal{S} = \mathcal{S}'\) .

让 \(j_{i}\) 为 \(B_{i}\) 进入。 \(A_{i}\) 的积，设j为B到E的典范单射，设 \(p_{i}\) 为 E 在 \(A_{i}\) ，并且让 \(p_{i}'\) 为 B 在 \(B_{i}\) 上的投影。于是我们有 \(p_{i} \circ j = j_{i} \circ p_{i}'\) 对于所有 \(i \in I\) 。根据CST10， \(\mathcal{G}\) 是关于该族结构的初始结构 \((A_{i}, \mathcal{G}_{i}, p_{i} \circ j)_{i \in I}\) ，以及 \(\mathcal{G}'\) 是关于该族结构的初始结构 \((A_{i}, \mathcal{G}_{i}, j_{i} \circ p_{i}')_{i \in I}\) 。因此结论成立。

逆像与积的概念通过以下判据相联系：

CST15. 设 \((\mathbf{A}_i)_{i\in I}\) 是一个集合族，并且对每个 \(i\in I\) 让 \(\mathcal{S}_i\) 是物种 \(\Sigma\) 在 \(\mathbf{A}_i\) 上的一个结构，并且让 \(f_i\) 为集合 E 到 \(\mathbf{A}_i\) 的一个映射。假设在积集 \(\mathbf{A} = \prod_{i\in I}\mathbf{A}_i\) 上存在一个积结构 \(\mathcal{S}\) ，它是

族 \((\mathcal{S}_t)_{t\in I}\) 的积结构。那么，关于族 \((A_t,\mathcal{S}_t,f_t)_{t\in I}\) 存在一个初始结构，当且仅当结构 \(\mathcal{S}\) 在映射 \(x\to f(x) = (f_t(x))\) （从 E 到 A）下具有逆像，并且此时这两个结构相同。

由于 \(f_{t} = \mathrm{pr}_{t} \circ f\) ，该判据是CST10的一个特例。

备注。设 \((\mathcal{S}_{\lambda})_{\lambda \in L}\) 是物种的一族结构 \(\Sigma\) 在同一集合A上；令 \(A_{\lambda}\) 表示集合A赋予结构 \(\mathcal{S}_{\lambda}\) ，并且让 \(I_{\lambda}\) 表示A到 \(A_{\lambda}\) 的恒等映射。设B为乘积集 \(A^{L} = \prod_{\lambda \in L} A_{\lambda}\) ，并且让 \(\Delta\) 为该乘积的对角线（第二章，

§ 5，第 3 号）。设 h 为 A 到 \(\Delta\) ，因此 \(h(x)\) 的元素 \((x_{\lambda})_{\lambda\in\mathbf{L}}\) 使得 \(x_{\lambda}=x\) 对于所有 \(\lambda\in L\) 。假设在B上存在一个乘积结构 \(G'\) 该族中 \((\mathcal{G}_{\lambda})\) 由于 h 是单射的，准则 CST 15 表明存在一个初始结构 G ，关于该族 \((A_{\lambda},\mathcal{G}_{\lambda},I_{\lambda})_{\lambda\in\mathbf{L}}\) ，当且仅当存在一个结构 \(G''\) 在 \(\Delta\) 上、由 \(G'\) 所诱导； \(G''\) 则与通过 h 将 G 传输到 \(\Delta\) 所得到的结构相同。特别地，当所有结构 \(G_{\lambda}\) 是相同的，则h是A（赋予此结构）到 \(\Delta\) 的一个同构 .

我们还有以下判据：

CST16. 设 \((\mathbf{A}_t)_{t \in \mathbf{I}}\) , \((\mathbf{B}_t)_{t \in \mathbf{I}}\) 为两个由同一集合指标化的集族。对每个 \(t \in \mathbf{I}\) ，让 \(\mathcal{S}_t\) 是物种 \(\Sigma\) 在 \(\mathbf{A}_t\) 上的一个结构，并令 \(\mathcal{S}_t'\) 是物种 \(\Sigma\) 在 \(\mathbf{B}_t\) 上的一个结构。假设在 \(\mathbf{A} = \prod_{t \in \mathbf{I}} \mathbf{A}_t\) （分别地，在 \(\mathbf{B} = \prod_{t \in \mathbf{I}} \mathbf{B}_t\) ) 一个积结构 \(\mathcal{S}\) （分别） \(\mathcal{S}'\) ) 该族的 \((\mathcal{S}_t)\) （分别） \((\mathcal{S}_t')\) ). 对于每一个 \(t \in \mathbf{I}\) ，让 \(f_t\) 为 \(\mathbf{A}_t\) 进入。 \(\mathbf{B}_t\) 。则映射 \(f = (f_t)_{t \in \mathbf{I}}\) 是……的一个态射 \(\mathbf{A}\) 进入。 \(\mathbf{B}\) .

让 \(p_{i}\) （分别） \(q_{i}\) ) 为 A 在 \(A_{i}\) （相应地，B 在 \(B_{i}\) ）上的投影。于是我们有 \(q_{i} \circ f = f_{i} \circ p_{i}\) 。由于 \(f_{i}\) 并且 \(p_{i}\) 是态射（准则

CST9）， \(f_{\iota} \circ p_{\iota}\) 是一个态射，通过 \((\mathbf{MO}_{\mathrm{II}})\) ；因此 \(f\) 由 (IN) 知是态射。

注记。对于大多数常见结构，CST 16 中给出的条件不仅是充分的，而且对于 \(f\) 成为态射（参见练习7）。特别地，在以下情形中确实如此（这些情形会出现，例如，如果 \(\Sigma\) 是序结构的种，* 或群结构的种，或拓扑结构的种 *，等等；参见练习8）：

存在一个族 \((a_{\iota})_{\iota\in I}\) 使得 \(a_{\iota}\in A_{\iota}\) 对于所有 \(\iota \in I\) 并且使得，如果我们令 \(r_\iota(x_\iota) = (y_\kappa)\) ，其中 \(y_{\iota} = x_{\iota}\) 并且 \(y_{\kappa} = a_{\kappa}\) 每当 \(\kappa\neq \iota\) ，每个映射 \(r_\iota\) 是……的一个态射 \(A_{\iota}\) 进入。 \(A\) 。因为如果 \(f = (f_\iota)\) 是……的一个态射 \(A\) 进入。 \(B\) ，我们可以写 \(f_{\iota} = q_{\iota}\circ f\circ r_{\iota}\) 对于所有 \(\iota \in I\) ，并且只需应用 \((MO_{II})\) .

注意， \(r_t\) 是态射，如果满足以下条件：

\begin{enumerate}
\def\labelenumi{(\alph{enumi})}
\tightlist
\item
  对于每个配备了种 \(\Sigma\) 的结构的集合 E，常值映射 \(z \to a_{t}\) 的态射，这完成了证明。 \(\mathbf{A}_{t}\) ；即，对于每个 \(\kappa \in I\) , \(p_{\kappa} \circ r_{t}\) 是……的一个态射 \(\mathbf{A}_{t}\) 进入。 \(\mathbf{A}_{\kappa}\) ，因为当 \(\kappa = t\) ，以及一个常值映射 \(z \to a_{\kappa}\) 当 \(\kappa \neq t\) ；根据积结构的定义， \(r_{t}\) 因此是 \(\mathbf{A}_{t}\) 到A的一个态射。
\end{enumerate}

上述列出的例子不仅满足(a)，还满足以下条件：

\begin{enumerate}
\def\labelenumi{(\alph{enumi})}
\setcounter{enumi}{1}
\tightlist
\item
  在每个集合 \(\mathbf{A}_i' = \mathbf{A}_i \times \prod_{x \neq i} \{a_x\}\) 上，结构 \(\mathcal{S}\) 诱导出一个物种结构 \(\Sigma\) .
\end{enumerate}

设 \(p_{i}^{\prime}\) 表示 \(p_{i}\) 在 \(A_{i}^{\prime}\) 上的限制。如果条件（a）和（b）都满足，则 \(p_{i}^{\prime}\) 是从 \(A_{i}^{\prime}\) 到 \(A_{i}\) 上的同构。事实上， \(p_{i}^{\prime} = p_{i} \circ j_{i}\) ，其中 \(j_{i}\) 是从 \(A_{i}^{\prime}\) 到 \(A\) 的典范注入， \(p_{i}^{\prime}\) 因而由 \((\mathrm{MO}_{\mathrm{II}})\) 可知是态射。另一方面， \(r_{i} = j_{i} \circ (p_{i}^{\prime})^{-1}\) ；因此 \((p_{i}^{\prime})^{-1}\) 是从 \(A_{i}\) 到 \(A_{i}^{\prime}\) 的态射，这是诱导结构定义的直接结论。

最后，我们有以下判据，它在许多情形下刻画了态射：

CST17. 设A和B是两个集合，赋予结构 \(\mathcal{S}_{\mathbf{A}}\) , \(\mathcal{S}_{\mathbf{B}}\) 同一类型的结构 \(\Sigma\) 。假设在 \(\mathbf{A} \times \mathbf{B}\) 该结构 \(\mathcal{S}_{\mathbf{A} \times \mathbf{B}}\) ，即 \(\mathcal{S}_{\mathbf{A}}\) 并且 \(\mathcal{S}_{\mathbf{B}}\) . 设 \(f\) 是A到B的一个映射，设F为其图像，并设 \(\pi\) 为双射 \(x \to (x, f(x))\) 将A映射到F上。然后，对于 \(f\) 要成为从A到B的态射，必要且充分条件是F上应存在一种物种结构。 \(\Sigma\) 由...诱导的 \(\mathcal{S}_{\mathbf{A} \times \mathbf{B}}\) 并且，当F被赋予这个结构时， \(\pi\) 应当是A到F的一个同构。

为了证明充分性，设 \(j\) 为 \(\mathbf{F}\) 进入。 \(\mathbf{A} \times \mathbf{B}\) 的典范嵌入。我们可以写成 \(f = \operatorname{pr}_2 \circ j \circ \pi\) ，以及 \(f\) 根据假设，则是由三个态射复合而成。

为证明必要性，令 \(\mathcal{G}_{\mathbf{F}}\) 为种类 \(\Sigma\) 的结构，它是利用双射将结构 \(\mathcal{G}_{\mathbf{A}}\) 传递到 \(\mathbf{F}\) 上所得，该双射为 \(\pi\) (§1，第5号)。我们必须证明 \(\mathcal{G}_{\mathbf{F}}\) 是 \(\mathcal{G}_{\mathbf{A} \times \mathbf{B}}\) 在 \(\mathbf{F}\) 上诱导的结构。首先注意， \(j\) 是从 \(\mathbf{F}\) 到 \(\mathbf{A} \times \mathbf{B}\) 的态射；事实上， \(j \circ \pi\) 是映射 \(x \to (x, f(x))\) ，它从 \(\mathbf{A}\) 到 \(\mathbf{A} \times \mathbf{B}\) ，并且根据对 \(f\) 的假设及乘积结构的定义是态射；因而由结构的定义可知， \(\mathcal{G}_{\mathbf{F}}, j\) 是态射。还需证明：若 \(\mathbf{E}\) 是赋予种类 \(\Sigma\) 结构的集合， \(g\) 是从 \(\mathbf{E}\) 到 \(\mathbf{F}\) 的映射，且 \(j \circ g\) 是从 \(\mathbf{E}\) 到 \(\mathbf{A} \times \mathbf{B}\) 的态射，则 \(g\) 也是态射；等价地， \(g_{1} = \pi \circ g\) 应是从 \(\mathbf{E}\) 到 \(\mathbf{A}\) 的态射。但 \(g_{1} = \mathrm{pr}_{1} \circ (j \circ g)\) ，因此结论由假设和乘积结构的定义立即得出。

